\textcolor{black}{\section{Introducción}}

\noindent Las tierras raras están conformadas por 17 elementos de la tabla periódica. En el último tercio del siglo presente, se ha logrado obtener cantidades suficientes de tierras raras con elevada pureza que ha dado lugar a materiales avanzados o tecnologías presentes hoy en la vida cotidiana. Algunos ejemplos son, materiales ópticos (tales como las fibras, los pogmentos, etc.), pantallas de televisión, imanes permanentes y catalizadores\cite{Cheng2018}. 
Las propiedades que hacen que las tierras raras sean de interés están íntimamente relacionadas con su configuración electrónica. Sus electones de interés óptico se encuentran principalmente en la subcapa \textit{4f}, la configuración electrónica general del átomo neutro es de la forma: \([Xe]\,4f^n6s^2\). Cuando se habla de tierras raras, normalmente nos interesa hablar de la configuración electrónica del ión dopante, como omo \(Eu^{3+}\), \(Dy^{3+}\), \(Tb^{3+}\) o \(Nd^{3+}\). 
Ahora bien, los electrones \textit{4f} son responsables de algunas de las líneas de emisión que producen las tierras raras y los electrones de dichas capa están protegidos por capas más externas \textit{5s} y \textit{5p}\cite{Cheng2018}. Si los electrones \textit{4f} estuvieran totalmente expuestos al entorno en donde se introducen los iones dopantes, el campo eléctrico del entorno podría modificar las energías que se observan. Debido a esta protección, los niveles de energía de los niveles \textit{4f} suelen conservar características similares al ser introducidos a materiales distintos.
\subsection{Espectro de emisión}
Al incidir luz sobre un ión de tierra rara, el ión alcanza un estado de mayor energía y después regresa a un estado de menor energía, en este proceso puede ocurrir dos posibilidades; por relajación radiactiva, el electrón pasa a un estado de menor energía y emite un fotón con una energía igual a la diferencia de energía entre niveles, pot relajación no radiativa, el ion pierde energía que se transmite en forma de vibraciones al material o a otro ion. En la relajación no radiativa destacan la relajación multifonónica, la relajación cruzada y la conversión ascendente cooperativa\cite{RighiniFerrari2005}. Estos mecanismos pueden competir con la emisión luminosa y modificar tanto la intensidad de fluorescencia como la eficiencia cuántica del material.
En un espectro de emisión se registran las longitudes emitidas después de la exitación cuando se la relación radiativa, la asignación de las bandas a transiciones electrónica permite reconocer la respuesta característica de cada ion.
\subsection{Fluorescencia}
La fluorescencia es la emisión de luz de un material, que ocurre cuando un material absorbe energía y libera parte de en forma de luz. Se distingue de la fosforescencia por la rapidez en que ocurre la emisión una vez retirada la excitación. \cite{RighiniFerrari2005}
\subsection{Obejtivos}
En este trabajo se caracterizan las propiedades de emisión de los iones
Dy$^{3+}$ y Sm$^{3+}$ mediante el análisis de sus espectros de
fotoluminiscencia. Se identifican las principales bandas de emisión,
se determinan sus longitudes de onda y números de onda, y se relacionan
con las transiciones electrónicas correspondientes. Además, se compara
la emisión del Dy$^{3+}$ en dos matrices vítreas distintas para evaluar
la influencia del material anfitrión, y se contrastan las mediciones
del Sm$^{3+}$ obtenidas con dos aparatos diferentes para examinar la
consistencia de las posiciones de las bandas registradas.